\section{Regular expander variants}
\label{sec:regular}

Sections~\ref{sec:binary-ea} and~\ref{sec:binary-ec} separate the distance
proof into sparse, intermediate, and dense message weights.  The dense regime
already attains the GV exponent.  The remaining expander cost comes from small
supports.  Region-stratified regularity reduces the chance that their edges
cluster in a small part of the output.

\subsection{The regional shell law}

Assume $n=d_L\ell$ and sample the left-regular expander from
Section~\ref{sec:constructions}.  Fix a binary message of weight $r$.  Within
one region, its image is the parity of $r$ independent uniform unit vectors in
$\F_2^\ell$.  A unit vector has one coordinate equal to one and all other
coordinates equal to zero.

Let $\pi^{(\ell)}_{r,u}$ be the probability that this regional image has
weight $u$.  Initialize $\pi^{(\ell)}_{0,0}:=1$.  Terms outside
$\{0,\ldots,\ell\}$ are zero, and
\begin{equation}
 \pi^{(\ell)}_{r+1,v}=
 \pi^{(\ell)}_{r,v-1}\frac{\ell-v+1}{\ell}
 +\pi^{(\ell)}_{r,v+1}\frac{v+1}{\ell}.
 \label{eq:regular-shell-recurrence}
\end{equation}

\begin{lemma}[Regional shell law]
\label{lem:regular-shell-law}
For a fixed weight-$r$ message, the $d_L$ regional images are independent.
Their weights have distribution $\pi^{(\ell)}_{r,\cdot}$.  Conditional on
weight $u$, a regional image is uniform among the binary vectors of weight
$u$.
\end{lemma}

\begin{proof}
If the current regional image has weight $u$, the next unit vector decreases
its weight with probability $u/\ell$ and increases it otherwise.  This proves
\eqref{eq:regular-shell-recurrence}.  Coordinate-permutation symmetry gives
conditional uniformity.  Incidences in distinct regions are independent.
\end{proof}

The recurrence is exact and contains only nonnegative quantities.  It is
practical for small message supports.  Two comparison bounds cover larger
supports.  Define
\begin{equation}
 a^{\Reg}_r:=\frac{1-e^{-2r/\ell}}2,
 \qquad
 \gamma_r:=\frac{e^r r!}{r^r}.
 \label{eq:regular-poisson-parameters}
\end{equation}

\begin{lemma}[Poisson comparison]
\label{lem:regular-poisson}
Let $\mathbf Z\in\F_2^n$ have independent $\Ber(a^{\Reg}_r)$ coordinates.
For every nonnegative function $F:\F_2^n\to\mathbb R_{\geq0}$,
\begin{equation}
 \E_{\B}[F(\x\B)]
 \leq\gamma_r^{d_L}\E[F(\mathbf Z)].
 \label{eq:regular-poisson}
\end{equation}
\end{lemma}

\begin{proof}
In each region, replace the $r$ unit-vector samples by independent Poisson
counts of mean $r/\ell$.  Conditional on total count $r$, the counts have the
original multinomial law.  Without conditioning, their parities are
independent $\Ber(a^{\Reg}_r)$ variables.  A Poisson variable of mean $r$
equals $r$ with probability $\gamma_r^{-1}$.  Dropping the independent
conditioning event in all $d_L$ regions proves the claim.
\end{proof}

Stirling's formula gives $\gamma_r=\Theta(\sqrt r)$.  Thus the comparison
cost is polynomial in $r$ when $d_L$ is fixed.

The next lemma bounds the probability of each particular expanded vector,
rather than only its weight.  We call such a probability a point mass.

\begin{lemma}[Regular point-mass bound]
\label{lem:regular-point-mass}
For every fixed weight-$r$ message and every $\y\in\F_2^n$,
\begin{equation}
 \Pr_{\B}[\x\B=\y]
 \leq 2^{d_L-n}\left(1+e^{-2r/\ell}\right)^n.
 \label{eq:regular-point-mass}
\end{equation}
\end{lemma}

\begin{proof}
For one region, test a vector by the parity of a fixed set of $j$
coordinates.  One uniform unit step preserves this parity minus the
probability that it flips it, giving average $1-2j/\ell$; after $r$
independent steps the average is $(1-2j/\ell)^r$.  Fourier inversion and
$1-x\leq e^{-x}$ bound every regional point probability by
$2^{1-\ell}(1+e^{-2r/\ell})^\ell$.  Multiply the bounds for the independent
regions.
\end{proof}

\subsection{Application to binary EA}

The accumulator state must cross regional boundaries.  Let
\begin{equation}
 \mathbf U(X,Y):=
 \begin{pmatrix}1&XY\\X&Y\end{pmatrix}.
 \label{eq:regular-ea-input-transfer}
\end{equation}
For a uniform regional input of weight $u$, define
\begin{align}
 \mathbf S^{\EA}_{\ell,u}(Y)
 &:=\binom{\ell}{u}^{-1}[X^u]\mathbf U(X,Y)^\ell,
 \label{eq:regular-ea-slice}\\
 \mathbf R^{\EA}_{\ell,r}(Y)
 &:=\sum_{u=0}^{\ell}\pi^{(\ell)}_{r,u}
 \mathbf S^{\EA}_{\ell,u}(Y).
 \label{eq:regular-ea-region}
\end{align}

\begin{theorem}[Exact left-regular EA first moment]
\label{thm:regular-ea-first-moment}
For $\G=\B\A$ with a region-stratified left-regular expander,
\begin{equation}
 \Pr_{\B}[\Bad_L(\G)]\leq
 \sum_{r=1}^{k}\binom{k}{r}\sum_{h=0}^{L}[Y^h]\,
 \e_0^{\mathsf T}
 \mathbf R^{\EA}_{\ell,r}(Y)^{d_L}\one.
 \label{eq:regular-ea-first-moment}
\end{equation}
\end{theorem}

\begin{proof}
Lemma~\ref{lem:regular-shell-law} expresses one regional input as a mixture of
uniform Hamming shells.  Matrix \eqref{eq:regular-ea-slice} averages the
accumulator transfer over one shell.  Matrix multiplication preserves the
accumulator state between regions.  Sum the low-weight coefficients and apply
\eqref{eq:first-moment}.
\end{proof}

For intermediate supports, Lemma~\ref{lem:regular-poisson} bounds the output
moment by
\begin{equation}
 \gamma_r^{d_L}\e_0^{\mathsf T}
 \mathbf T_{a^{\Reg}_r}(z)^n\one.
 \label{eq:regular-ea-poisson-moment}
\end{equation}
For dense supports, invertibility of the accumulator and
Lemma~\ref{lem:regular-point-mass} give
\begin{equation}
 \Pr[\wt(\x\B\A)\leq L]
 \leq V_{2,n}(L)2^{d_L-n}(1+e^{-2r/\ell})^n,
 \quad
 V_{2,n}(L):=\sum_{h=0}^{L}\binom nh.
 \label{eq:regular-ea-dense}
\end{equation}

\subsection{Application to binary EC}

For a regional convolution input of weight $u$, define
\begin{align}
 \mathbf S^{\EC}_{\ell,u}(Y)
 &:=\binom{\ell}{u}^{-1}[X^u]\mathbf W_m(X,Y)^\ell,
 \label{eq:regular-ec-slice}\\
 \mathbf R^{\EC}_{\ell,r}(Y)
 &:=\sum_{u=0}^{\ell}\pi^{(\ell)}_{r,u}
 \mathbf S^{\EC}_{\ell,u}(Y).
 \label{eq:regular-ec-region}
\end{align}

\begin{theorem}[Exact left-regular wrapped-EC first moment]
\label{thm:regular-ec-first-moment}
Independently sample a region-stratified left-regular expander $\B$ and a
wrapped convolution $\C$ of memory $m$.  Then
\begin{equation}
 \Pr_{\B,\C}[\Bad_L(\B\C)]\leq
 \sum_{r=1}^{k}\binom{k}{r}\sum_{h=0}^{L}[Y^h]\,
 \e_m^{\mathsf T}
 \mathbf R^{\EC}_{\ell,r}(Y)^{d_L}\one.
 \label{eq:regular-ec-first-moment}
\end{equation}
\end{theorem}

\begin{proof}
The proof is the regional version of
Theorem~\ref{thm:ec-exact-first-moment}.  Lemma~\ref{lem:regular-shell-law}
gives the input mixture in each region.  Matrix multiplication retains the
convolution state at regional boundaries.  The first-moment bound completes
the proof.
\end{proof}

The Poisson and dense bounds become
\begin{align}
 \E[z^{\wt(\x\B\C)}]
 &\leq\gamma_r^{d_L}
 \e_m^{\mathsf T}\mathbf N_{m,a^{\Reg}_r}(z)^n\one,
 \label{eq:regular-ec-poisson}\\
 \Pr[\wt(\x\B\C)\leq L]
 &\leq V_{2,n}(L)2^{d_L-n}(1+e^{-2r/\ell})^n.
 \label{eq:regular-ec-dense}
\end{align}
Equation~\eqref{eq:regular-ec-dense} holds for every fixed invertible binary
map $\C$ and therefore also after averaging the convolution.

\subsection{Two-sided regular incidence over the binary field}
\label{sec:binary-two-sided-regular}

Fix a right degree $d_R$ and let $k=d_R\ell$.  In each region, the
two-sided regular expander partitions the $k$ left vertices into $\ell$
groups of size $d_R$.  A right coordinate is one when its group contains an
odd number of vertices from the message support.  Define
\begin{align}
 E_{d_R}(X)&:=\frac{(1+X)^{d_R}+(1-X)^{d_R}}2,
 &
 O_{d_R}(X)&:=\frac{(1+X)^{d_R}-(1-X)^{d_R}}2.
 \label{eq:binary-biregular-parity-polynomials}
\end{align}
The coefficient of $X^j$ in $E_{d_R}$ is $\binom{d_R}{j}$ for even $j$
and zero for odd $j$.  The polynomial $O_{d_R}$ gives the complementary
odd coefficients.

Write the accumulator transfer as
$\mathbf U(X,Y)=\mathbf U_0(Y)+X\mathbf U_1(Y)$.  Similarly, write the
wrapped-convolution transfer as
$\mathbf W_m(X,Y)=\mathbf W_{m,0}(Y)+X\mathbf W_{m,1}(Y)$.  For
$\mathsf R\in\{\EA,\EC\}$, let $(\mathbf Q_0,\mathbf Q_1)$ denote the
corresponding pair of input-conditioned matrices.  Define
\begin{equation}
 \mathbf R^{\mathsf R,\mathrm{bi}}_{\ell,d_R,r}(Y):=
 \binom{k}{r}^{-1}[X^r]
 \left(E_{d_R}(X)\mathbf Q_0(Y)
       +O_{d_R}(X)\mathbf Q_1(Y)\right)^\ell.
 \label{eq:binary-biregular-region-transfer}
\end{equation}

\begin{theorem}[Exact two-sided regular binary first moment]
\label{thm:binary-biregular-first-moment}
Let $n=d_L\ell$.  Sample $\B$ from the two-sided regular ensemble by
sampling the regional permutations independently.  For binary EA,
\begin{equation}
 \Pr_{\B}[\Bad_L(\B\A)]\leq
 \sum_{r=1}^{k}\binom{k}{r}\sum_{h=0}^{L}[Y^h]\,
 \e_0^{\mathsf T}
 \mathbf R^{\EA,\mathrm{bi}}_{\ell,d_R,r}(Y)^{d_L}\one.
 \label{eq:binary-biregular-ea-first-moment}
\end{equation}
For binary wrapped EC, independently sample a wrapped convolution $\C$ of
memory $m$ and replace $\EA$ by
$\EC$ in \eqref{eq:binary-biregular-region-transfer}.  Then
\begin{equation}
 \Pr_{\B,\C}[\Bad_L(\B\C)]\leq
 \sum_{r=1}^{k}\binom{k}{r}\sum_{h=0}^{L}[Y^h]\,
 \e_m^{\mathsf T}
 \mathbf R^{\EC,\mathrm{bi}}_{\ell,d_R,r}(Y)^{d_L}\one.
 \label{eq:binary-biregular-ec-first-moment}
\end{equation}
\end{theorem}

\begin{proof}
Fix a message support of size $r$.  Its positions in one regional
permutation form a uniform $r$-subset of the $k$ slots.  Within one group,
$E_{d_R}$ counts selections that produce input zero, and $O_{d_R}$ counts
selections that produce input one.  The coefficient of $X^r$ in
\eqref{eq:binary-biregular-region-transfer} counts all regional selections
and retains the state of the recursive map.  Division by $\binom{k}{r}$
averages over the uniform subset.  The regional permutations are independent,
so matrix multiplication composes their transfers.  Summing over message
supports and low output weights gives the two bounds.
\end{proof}

The all-one message imposes a parity restriction that has no labeled
large-field analogue.

\begin{proposition}[Binary right-degree parity]
\label{prop:binary-right-degree-parity}
For a binary two-sided regular expander,
\begin{equation}
 \one\B=(d_R\bmod 2)\one.
 \label{eq:binary-right-degree-parity}
\end{equation}
If $d_R$ is even, both $\B\A$ and $\B\C$ are noninjective and have minimum
distance zero.
\end{proposition}

\begin{proof}
Every right coordinate has exactly $d_R$ incident edges.  This proves
\eqref{eq:binary-right-degree-parity}.  If $d_R$ is even, the nonzero
all-one message maps to zero before the invertible recursive map.
\end{proof}

At rate one half, $d_L=2d_R$.  A usable binary two-sided regular ensemble
therefore requires $d_L\equiv2\pmod4$.  This excludes the current left-regular
degrees $64$ for EA and $28$ for EC.

For right degree five, a local concentration argument removes the
square-root loss from the positive coefficient saddle.  Fix
$r\in\{1,\ldots,k-1\}$ and put $\theta_r:=r/k$.  Let
\begin{equation}
 q_1(\theta):=\frac{1-(1-2\theta)^5}{2},
 \qquad q_0(\theta):=1-q_1(\theta).
 \label{eq:degree-five-parity-probabilities}
\end{equation}
For $b\in\{0,1\}$, define a random variable $Z_b\in\{0,1,2\}$ by
\begin{equation}
 \Pr[Z_b=j]:=
 \frac{\binom{5}{2j+b}\theta_r^{2j+b}(1-\theta_r)^{5-2j-b}}
 {q_b(\theta_r)}.
 \label{eq:degree-five-conditioned-half-count}
\end{equation}
Thus $2Z_b+b$ is the number of selected slots in one group, conditioned on
parity $b$.  Define
\begin{align}
 v_*(\theta_r)&:=\min_{b\in\{0,1\}}\operatorname{Var}(Z_b),
&
 \beta_{k,r}&:=\binom{k}{r}\theta_r^r(1-\theta_r)^{k-r},
 \label{eq:degree-five-local-parameters}\\
 M_{k,r}&:=
 \frac{\max\{q_0(\theta_r),q_1(\theta_r)\}^{\ell}}{\beta_{k,r}}
 \min\left\{1,\sqrt{\frac{\pi}{8\ell v_*(\theta_r)}}\right\}.
 \label{eq:degree-five-point-mass}
\end{align}

\begin{lemma}[Point mass for right degree five]
\label{lem:degree-five-point-mass}
Let $\mathbf Y\in\F_2^\ell$ be the image of a fixed weight-$r$ support in
one region of a two-sided regular expander with $d_R=5$.  Then
\begin{equation}
 \max_{\y\in\F_2^\ell}\Pr[\mathbf Y=\y]\leq M_{k,r}.
 \label{eq:degree-five-regional-point-mass}
\end{equation}
\end{lemma}

\begin{proof}
Independently select each of the $k$ regional slots with probability $\theta_r$.
Conditioning on exactly $r$ selected slots gives the uniform $r$-subset in
the two-sided regular ensemble.  Before conditioning, the group parities are
independent and have probabilities $q_0(\theta_r)$ and $q_1(\theta_r)$.

Fix a parity vector $\y$ of weight $u$.  Conditional on $\mathbf Y=\y$, the
selected-slot count equals
\begin{equation}
 u+2\sum_{i=1}^{\ell} Z_{y_i}.
 \label{eq:degree-five-conditioned-total}
\end{equation}
The probability-generating polynomials of $Z_0$ and $Z_1$, before
normalization, are
\begin{align*}
 &(1-\theta_r)^5+10\theta_r^2(1-\theta_r)^3T
   +5\theta_r^4(1-\theta_r)T^2,\\
 &5\theta_r(1-\theta_r)^4+10\theta_r^3(1-\theta_r)^2T
   +\theta_r^5T^2.
\end{align*}
Their discriminants are $80\theta_r^4(1-\theta_r)^6$ and
$80\theta_r^6(1-\theta_r)^4$.  Thus both polynomials have two real negative
roots.
After normalization, each $Z_b$ has the distribution of a sum of two
independent Bernoulli variables.

Consequently, the sum in \eqref{eq:degree-five-conditioned-total} is a sum
of independent Bernoulli variables with variance at least
$\ell v_*(\theta_r)$.  If such a sum has variance $V$, Fourier inversion gives
\begin{equation}
 \max_j\Pr[Z=j]
 \leq\frac1{2\pi}\int_{-\pi}^{\pi}
 e^{-2V\sin^2(t/2)}\,dt
 \leq\sqrt{\frac{\pi}{8V}}.
 \label{eq:poisson-binomial-local-mass}
\end{equation}
The last inequality uses $\sin(|t|/2)\geq |t|/\pi$ on this interval.
Bayes' rule now gives
\begin{equation*}
 \Pr[\mathbf Y=\y\mid\text{$r$ selected slots}]
 \leq
 \frac{\max\{q_0(\theta_r),q_1(\theta_r)\}^{\ell}}{\beta_{k,r}}
 \min\left\{1,\sqrt{\frac{\pi}{8\ell v_*(\theta_r)}}\right\}.
\end{equation*}
This is \eqref{eq:degree-five-regional-point-mass}.
\end{proof}

The certificate evaluates \eqref{eq:degree-five-regional-point-mass} in
blocks.  The following endpoint bound makes that evaluation rigorous.  Put
$r_{\mathrm{mid}}:=(k-1)/2$ and $\theta_a:=a/k$.  For
$k/5\leq a\leq b\leq r_{\mathrm{mid}}$, define
\begin{equation}
 \widehat M_{a,b}:=
 \frac{q_0(\theta_a)^\ell}{\beta_{k,b}}
 \sqrt{\frac{\pi}{8\ell\operatorname{Var}(Z_1(\theta_a))}}.
 \label{eq:degree-five-block-point-mass}
\end{equation}
Here $Z_1(\theta)$ denotes the variable from
\eqref{eq:degree-five-conditioned-half-count} with parameter $\theta$.

\begin{lemma}[Endpoint bound for the central range]
\label{lem:degree-five-central-block}
For $k/5\leq a\leq b\leq r_{\mathrm{mid}}$ and every fixed invertible
$\mathbf D$,
\begin{equation}
 \sum_{r=a}^{b}\binom{k}{r}
 \Pr_{\B}[\wt(\x_r\B\mathbf D)\leq L]
 \leq
 (b-a+1)\binom{k}{b}V_{2,n}(L)\widehat M_{a,b}^{d_L},
 \label{eq:degree-five-central-block}
\end{equation}
where $\x_r$ is any fixed weight-$r$ message.  The same bound applies to
the complementary supports $k-b\leq r\leq k-a$.
\end{lemma}

\begin{proof}
For $\theta\in[1/5,1/2]$, $q_0(\theta)$ decreases.  In the same interval,
$v_*(\theta)=\operatorname{Var}(Z_1(\theta))$, and this variance increases.
To verify the latter claim, write
\begin{equation*}
 \operatorname{Var}(Z_1(\theta))=
 \frac{10\theta^2(1-\theta)^2
 (8\theta^4-24\theta^3+32\theta^2-20\theta+5)}
 {(16\theta^4-40\theta^3+40\theta^2-20\theta+5)^2}.
\end{equation*}
Its derivative has the same sign as
\begin{equation*}
 \theta(\theta-1)(2\theta-1)(4\theta^2-10\theta+5)
 (4\theta^4-20\theta^3+30\theta^2-20\theta+5),
\end{equation*}
 which is nonnegative on $[1/5,1/2]$.  Indeed, after setting $t=1-2\theta$,
 the last two factors become
 $t^2+3t+1$ and
 $(t^4+6t^3+6t^2+6t+1)/4$.
 The same substitution gives
 \begin{equation*}
  \operatorname{Var}(Z_0(\theta))-\operatorname{Var}(Z_1(\theta))=
  \frac{5t^3(1-t)^2(1+t)^2(t^2+1)(2t^4+t^2+2)}
  {4(t^4-t^3+t^2-t+1)^2(t^4+t^3+t^2+t+1)^2},
 \end{equation*}
 which is nonnegative for $t\in[0,3/5]$.  Hence
 $v_*(\theta)=\operatorname{Var}(Z_1(\theta))$ throughout the interval.

It remains to control $\beta_{k,r}$.  Its consecutive ratio is
\begin{equation*}
 \frac{\beta_{k,r+1}}{\beta_{k,r}}
 =\left(1+\frac1r\right)^r
  \left(1-\frac1{k-r}\right)^{k-r-1}\leq1
 \qquad(r\leq r_{\mathrm{mid}}).
\end{equation*}
 Here the second factor is the reciprocal of
 $(1+1/(k-r-1))^{k-r-1}$.  The inequality follows because
 $(1+1/s)^s$ increases with $s$ and $r\leq k-r-1$.
Consequently, every point mass in a block is at most
$\widehat M_{a,b}^{d_L}$.  The message count is at most $\binom{k}{b}$.
Lemma~\ref{lem:degree-five-point-mass}, invertibility of $\mathbf D$, and a
sum over the block prove \eqref{eq:degree-five-central-block}.  Replacing a
selected set by its complement exchanges even and odd group parities.  All
terms in the bound are therefore unchanged.
\end{proof}

\begin{corollary}[Dense degree-five bound]
\label{cor:degree-five-dense-bound}
Let $\mathbf D$ be any fixed invertible binary map on $\F_2^n$.  For a fixed
weight-$r$ message $\x$,
\begin{equation}
 \Pr_{\B}[\wt(\x\B\mathbf D)\leq L]
 \leq V_{2,n}(L)M_{k,r}^{d_L}.
 \label{eq:degree-five-dense-bound}
\end{equation}
The same bound holds when $\mathbf D$ is sampled independently of $\B$.
\end{corollary}

\begin{proof}
The regional permutations are independent, so every point probability of
$\x\B$ is at most $M_{k,r}^{d_L}$.  Invertibility gives exactly
$V_{2,n}(L)$ expanded vectors whose image under $\mathbf D$ has weight at most
$L$.  Summing their point probabilities proves the fixed-map claim.
Averaging over an independent $\mathbf D$ proves the second claim.
\end{proof}

Right degree three admits the same local argument with a simpler conditional
distribution.  For $\theta:=r/k$, define
\begin{equation}
 q^{(3)}_1(\theta):=\frac{1-(1-2\theta)^3}{2},
 \qquad q^{(3)}_0(\theta):=1-q^{(3)}_1(\theta).
 \label{eq:degree-three-parity-probabilities}
\end{equation}
For $s\in\{0,1\}$, let $Z^{(3)}_s(\theta)\in\{0,1\}$ satisfy
\begin{equation}
 \Pr[Z^{(3)}_s(\theta)=j]
 :=\frac{\binom{3}{2j+s}\theta^{2j+s}(1-\theta)^{3-2j-s}}
 {q^{(3)}_s(\theta)}.
 \label{eq:degree-three-conditioned-half-count}
\end{equation}
Thus $2Z^{(3)}_s(\theta)+s$ is the selected-slot count in a three-slot group,
conditioned on its parity.  Put
\begin{equation}
 M^{(3)}_{k,r}:=
 \frac{\max\{q^{(3)}_0(\theta),q^{(3)}_1(\theta)\}^{\ell}}{\beta_{k,r}}
 \min\left\{1,
 \sqrt{\frac{\pi}{8\ell\min_s\operatorname{Var}(Z^{(3)}_s(\theta))}}
 \right\}.
 \label{eq:degree-three-point-mass}
\end{equation}

\begin{lemma}[Central blocks for right degree three]
\label{lem:degree-three-central-block}
Let $r_{\mathrm{mid}}=(k-1)/2$, and let
$1\leq a\leq b\leq r_{\mathrm{mid}}$.  Suppose
\begin{equation*}
 \sqrt{\frac{\pi}
 {8\ell\operatorname{Var}(Z^{(3)}_1(a/k))}}\leq1.
\end{equation*}
Define
\begin{equation}
 \widehat M^{(3)}_{a,b}:=
 \frac{q^{(3)}_0(a/k)^\ell}{\beta_{k,b}}
 \sqrt{\frac{\pi}
 {8\ell\operatorname{Var}(Z^{(3)}_1(a/k))}}.
 \label{eq:degree-three-block-point-mass}
\end{equation}
Then, for every fixed invertible $\mathbf D$,
\begin{equation}
 \sum_{r=a}^{b}\binom{k}{r}
 \Pr_{\B}[\wt(\x_r\B\mathbf D)\leq L]
 \leq(b-a+1)\binom{k}{b}V_{2,n}(L)
 \bigl(\widehat M^{(3)}_{a,b}\bigr)^{d_L}.
 \label{eq:degree-three-central-block}
\end{equation}
The same bound applies to supports $k-b\leq r\leq k-a$.
\end{lemma}

\begin{proof}
Conditioning independent slot selections on their total count gives a uniform
weight-$r$ support.  Conditional on the parity vector, the remaining count is
a sum of the independent Bernoulli variables in
\eqref{eq:degree-three-conditioned-half-count}.  The Fourier bound
\eqref{eq:poisson-binomial-local-mass} therefore gives
\begin{equation*}
 \max_{\y\in\F_2^\ell}\Pr[\mathbf Y=\y]
 \leq M^{(3)}_{k,r}.
\end{equation*}

For $0<\theta\leq1/2$, direct calculation gives
\begin{align*}
 \operatorname{Var}(Z^{(3)}_0(\theta))
 &=\frac{3\theta^2(1-\theta)^2}{(4\theta^2-2\theta+1)^2},\\
 \operatorname{Var}(Z^{(3)}_1(\theta))
 &=\frac{3\theta^2(1-\theta)^2}{(4\theta^2-6\theta+3)^2}.
\end{align*}
The second variance is no larger than the first and increases on this
interval.  Indeed, their difference has the sign of $1-2\theta$, and
\begin{equation*}
 \frac{d}{d\theta}\operatorname{Var}(Z^{(3)}_1)
 =\frac{-6\theta(\theta-1)(2\theta^2-6\theta+3)}
 {(4\theta^2-6\theta+3)^3}\geq0.
\end{equation*}
Also, $q^{(3)}_0(\theta)$ decreases for $\theta\leq1/2$, and
$\beta_{k,r}$ decreases for integer $r\leq r_{\mathrm{mid}}$, as shown in
the proof of Lemma~\ref{lem:degree-five-central-block}.  Hence every regional point mass
in the block is at most $\widehat M^{(3)}_{a,b}$.  Independence of the
regions, invertibility of $\mathbf D$, and a sum over the block prove
\eqref{eq:degree-three-central-block}.  Complementing all selected slots
exchanges the two group parities and proves the final claim.
\end{proof}

The same argument extends to every odd right degree.  The following form is
weaker than the specialized degree-three and degree-five bounds, but it is
sufficient for the intermediate finite profiles.  Let $d=2h+1$ and define
\begin{align*}
 \Theta_0(d)&:=\left\{\frac{(2j+1)\pi}{2d}:0\leq j<h\right\},&
 \Theta_1(d)&:=\left\{\frac{j\pi}{d}:1\leq j\leq h\right\},\\
 f(t)&:=\frac{t}{(1+t)^2},&c_a&:=\frac{k-a}{a}.
\end{align*}
For $s\in\{0,1\}$, put
\begin{equation}
 v_{d,s}(a):=
 \sum_{\theta\in\Theta_s(d)}
 \min\left\{f\bigl(c_a^2\tan^2\theta\bigr),
                   f\bigl(\tan^2\theta\bigr)\right\},
 \qquad
 v_d(a):=\min_s v_{d,s}(a).
 \label{eq:odd-degree-uniform-variance}
\end{equation}

\begin{lemma}[Central blocks for odd right degree]
\label{lem:odd-degree-central-block}
Let $d_R=d$ be odd, let $r_{\mathrm{mid}}=(k-1)/2$, and let
$1\leq a\leq b\leq r_{\mathrm{mid}}$.  Define $\theta_a:=a/k$ and
\begin{equation}
 \widehat M^{(d)}_{a,b}:=
 \frac{q_{0,d}(\theta_a)^\ell}{\beta_{k,b}}
 \sqrt{\frac{\pi}{8\ell v_d(a)}},
 \qquad
 q_{0,d}(\theta):=\frac{1+(1-2\theta)^d}{2}.
 \label{eq:odd-degree-block-point-mass}
\end{equation}
Suppose the square-root factor in
\eqref{eq:odd-degree-block-point-mass} is at most one.  For every fixed
invertible $\mathbf D$,
\begin{equation}
 \sum_{r=a}^{b}\binom{k}{r}
 \Pr_{\B}[\wt(\x_r\B\mathbf D)\leq L]
 \leq(b-a+1)\binom{k}{b}V_{2,n}(L)
 \bigl(\widehat M^{(d)}_{a,b}\bigr)^{d_L}.
 \label{eq:odd-degree-central-block}
\end{equation}
The same bound applies to supports $k-b\leq r\leq k-a$.
\end{lemma}

\begin{proof}
Fix $\theta\in[\theta_a,1/2]$.  In one right group, let $J$ count independently
selected slots.  Conditional on $J\bmod2=s$, the probability-generating
polynomial of $(J-s)/2$ is proportional to
\begin{equation*}
 \sum_{j=0}^{h}\binom{d}{2j+s}
 \theta^{2j+s}(1-\theta)^{d-2j-s}T^j.
\end{equation*}
For $s=0$, its roots are
\begin{equation*}
 -\left(\frac{1-\theta}{\theta}\right)^2\tan^2\varphi,
 \qquad \varphi\in\Theta_0(d).
\end{equation*}
For $s=1$, the same formula holds with $\varphi\in\Theta_1(d)$.  These
identities follow by taking the even part of
$(1-\theta+\theta\sqrt T)^d$ and by
dividing its odd part by $\sqrt T$.

All roots are negative.  After normalization, each polynomial is therefore
the generating polynomial of a sum of independent Bernoulli variables.  A
root $-\rho$ contributes variance $f(\rho)$.  As $\theta$ runs from
$\theta_a$ to
$1/2$, each $\rho$ changes monotonically between
$c_a^2\tan^2\varphi$ and $\tan^2\varphi$.  The function $f$ increases on
$(0,1]$ and decreases on $[1,\infty)$.  Its minimum on this interval is at
an endpoint.  Hence either conditional half-count has variance at least
$v_d(a)$.

The Fourier estimate \eqref{eq:poisson-binomial-local-mass} now bounds the
conditional total-count point mass.  For $\theta\leq1/2$,
$q_{0,d}(\theta)$ is the larger parity probability and decreases with
$\theta$.  Also, $\beta_{k,r}$ decreases for
$r\leq r_{\mathrm{mid}}$.  The endpoint argument
from Lemma~\ref{lem:degree-five-central-block} proves
\eqref{eq:odd-degree-central-block}.  Complementing all slots exchanges the
two parity laws and proves the final claim.
\end{proof}

\subsection{Certified improvements at rate one half}

\begin{theorem}[Regular binary EA and EC certificates]
\label{thm:regular-binary-certificates}
The following claims hold.
\begin{enumerate}
  \item Let $k=1{,}048{,}576$, $d_L=64$, $\ell=32{,}768$,
  $n=d_L\ell=2k$, and $L=230{,}718$.  Sample $\B$ from the
  region-stratified left-regular ensemble.  Then
  \begin{equation}
   \Pr_{\B}[\Bad_L(\B\A)]
   \leq8.982345581\cdot10^{-7}<2^{-20.0864}.
   \label{eq:regular-ea-certificate}
  \end{equation}
  \item Let $k=1{,}048{,}572$, $d_L=28$, $\ell=74{,}898$,
  $n=d_L\ell=2k$, and $L=230{,}730$.  Sample $\B$ from the
  region-stratified left-regular ensemble and independently sample a wrapped
  convolution $\C$ of memory $9$.  Then
  \begin{equation}
   \Pr_{\B,\C}[\Bad_L(\B\C)]
   <2.036627\cdot10^{-7}<2^{-22.2273}.
   \label{eq:regular-ec-certificate}
  \end{equation}
  \item Let $k=1{,}048{,}575$, $d_L=10$, $d_R=5$,
  $\ell=209{,}715$, $n=2k$, and $L=230{,}729$.  Sample $\B$ from the
  two-sided regular ensemble and independently sample a wrapped convolution
  $\C$ of memory $15$.  Then
  \begin{equation}
   \Pr_{\B,\C}[\Bad_L(\B\C)]
   <1.6475000\cdot10^{-10}<2^{-32.4990}.
   \label{eq:biregular-binary-ec-certificate}
  \end{equation}
  \item Let $k=1{,}048{,}635$, $d_L=14$, $d_R=7$,
  $\ell=149{,}805$, $n=2k$, and $L=230{,}742$.  Sample $\B$ from the
  two-sided regular ensemble and independently sample a wrapped convolution
  $\C$ of memory $6$.  Then
  \begin{equation}
   \Pr_{\B,\C}[\Bad_L(\B\C)]
   <1.402384\cdot10^{-8}<2^{-26.0875}.
   \label{eq:biregular-binary-ec-degree-fourteen-certificate}
  \end{equation}
  \item Let $k=1{,}048{,}635$, $d_L=18$, $d_R=9$,
  $\ell=116{,}515$, $n=2k$, and $L=230{,}742$.  Sample $\B$ from the
  two-sided regular ensemble and independently sample a wrapped convolution
  $\C$ of memory $4$.  Then
  \begin{equation}
   \Pr_{\B,\C}[\Bad_L(\B\C)]
   <2.987318\cdot10^{-7}<2^{-21.6746}.
   \label{eq:biregular-binary-ec-degree-eighteen-certificate}
  \end{equation}
  \item Let $k=1{,}048{,}575$, $d_L=6$, $d_R=3$,
  $\ell=349{,}525$, $n=2k$, and $L=230{,}729$.  Sample $\B$ from the
  two-sided regular ensemble and independently sample a wrapped convolution
  $\C$ of memory $79$.  Then
  \begin{equation}
   \Pr_{\B,\C}[\Bad_L(\B\C)]
   <3.558444\cdot10^{-7}<2^{-21.4222}.
   \label{eq:biregular-binary-ec-degree-six-certificate}
  \end{equation}
\end{enumerate}
\end{theorem}

\begin{proof}
For EA, the verifier combines the exact regional transfer for small supports,
\eqref{eq:regular-ea-poisson-moment} for intermediate supports, and
\eqref{eq:regular-ea-dense} for dense supports.  For EC it applies the
corresponding three bounds from
\eqref{eq:regular-ec-first-moment}--\eqref{eq:regular-ec-dense}.  The support
intervals cover $[1,k]$.

For the degree-$10/5$ EC result, the verifier uses exact regional transfers for
supports $1$ through $16$.  It uses $111$ positive coefficient blocks for
the two outer ranges and $47$ applications of
Lemma~\ref{lem:degree-five-central-block} for the central range.  It checks
the all-one support directly.  Outward-rounded Arb arithmetic gives the
displayed values.  Section~\ref{sec:certificates} gives the verifier details.

For the degree-$14/7$ and degree-$18/9$ results, exact regional transfers
cover supports $1$ through $24$ and $1$ through $32$, respectively.  Positive
coefficient blocks cover the two outer ranges.  The central ranges use
Lemma~\ref{lem:odd-degree-central-block}.  A separate transfer checks each
all-one support.

For the degree-$6/3$ result, exact regional transfers cover supports $1$
through $383$.  A layered representation allows each reachable matrix entry
to use an independent power-of-two scale.  The representation uses only positive
operations and rounds every operation upward.  It therefore retains entries
below the ordinary binary64 range without adding a matrix-wide absolute
error.  The verifier uses $103$ coefficient blocks outside the central range
and $28$ applications of Lemma~\ref{lem:degree-three-central-block} inside
that range.  It checks the all-one support separately.
Section~\ref{sec:certificates} gives the resulting bounds.
\end{proof}

The EA cutoff is $0.9998823119$ times the binary rate-one-half GV distance.
A matched Bernoulli certificate requires expected degree $75$.  Left
regularity therefore removes $14.67\%$ of the expander edges.

The EC cutoff is $0.9999381317$ times the binary rate-one-half GV distance.
The comparable Bernoulli memory-$9$ profile has expected degree $42.2127$.
Replacing it by degree $28$ removes $33.67\%$ of the expander edges while
losing eleven units of output weight.

The two-sided EC cutoff is $0.9999309371$ times the binary rate-one-half GV
distance.  Relative to the left-regular degree-$28$ profile, two-sided
regularity and six additional memory words remove $64.29\%$ of the expander
edges.  The simple work proxy $d_L+m$ decreases from $37$ to $25$.

Table~\ref{tab:binary-recursive-map-frontier} gives the certified binary
frontier.  We record the accumulator as a memory-one recursive map because it
retains only $y_{t-1}$.  Over $\F_2$, this map is also the memory-one instance
of the wrapped convolution in \eqref{eq:wrapped-convolution}.  Its certified
expander is left regular; the four other rows use two-sided regular incidence.
A dash in the $d_R$ column means that the distribution does not fix the right
degree.
The degree-$14/7$, memory-$6$ profile minimizes the proxy $d_L+m$, with value
$20$.  Reducing the left degree to $10$ or $6$ requires substantially more
convolution memory.

\begin{table}[t]
\centering
\caption{Certified binary recursive-map frontier at rate one half.}
\label{tab:binary-recursive-map-frontier}
\begin{tabular}{@{}llrrrrr@{}}
\toprule
map & expander & $d_L$ & $d_R$ & memory $m$ & $d_L+m$ & failure bits \\
\midrule
EA & left regular & $64$ & -- & $1$  & $65$ & $20.0864$ \\
EC & two-sided regular & $18$ & $9$ & $4$  & $22$ & $21.6746$ \\
EC & two-sided regular & $14$ & $7$ & $6$  & $20$ & $26.0875$ \\
EC & two-sided regular & $10$ & $5$ & $15$ & $25$ & $32.4990$ \\
EC & two-sided regular & $6$  & $3$ & $79$ & $85$ & $21.4222$ \\
\bottomrule
\end{tabular}
\end{table}

At rate one half, the degrees satisfy $d_L=2d_R$.  The binary construction
also requires odd $d_R$: if $d_R$ is even, the all-one message contributes an
even number of ones to every right vertex and therefore encodes to zero.
Consequently, degree $8/4$ is not an intermediate candidate between $6/3$
and $10/5$.  The usable left degrees at rate one half are congruent to two
modulo four; Table~\ref{tab:binary-recursive-map-frontier} lists four such EC
profiles.

The EA outcome is different.  At the same cutoff, the $d_L/d_R=62/31$
diagnostic is dominated by message weight two and does not reach a $2^{-20}$
first moment.  Fixing the right degrees does not recover the two removed EA
edges.

The degree-$6/3$ calculation also illustrates a numerical issue in exact
enumerator proofs.  Some reachable transfer entries lie below the ordinary
binary64 range.  A global underflow allowance becomes useless after the union
bound multiplies it by the number of messages.  The layered verifier avoids
this loss by retaining each exponent band as a separate positive matrix.
The matrix products use relative upward rounding; Arb evaluates the remaining
scalar expressions.

\subsection{Scalar extension preserves distance but not field mixing}

The binary generators above can act on messages over an extension field of
$\F_2$, and this reinterpretation preserves minimum distance.  It does not,
however, turn a binary encoder into a field-mixing encoder.

\begin{lemma}[Distance under scalar extension]
\label{lem:binary-scalar-extension}
Let $\G\in\F_2^{k\times n}$, and let $K$ be a finite extension of $\F_2$.
View the entries of $\G$ as elements of $K$.  Then
\begin{equation}
 \min_{\x\in K^k\setminus\{0\}}\wt(\x\G)
 =
 \min_{\x\in\F_2^k\setminus\{0\}}\wt(\x\G).
 \label{eq:binary-scalar-extension-distance}
\end{equation}
\end{lemma}

\begin{proof}
Fix an $\F_2$-basis $\beta_1,\ldots,\beta_t$ of $K$.  Every
$\x\in K^k$ has a unique decomposition
\begin{equation*}
 \x=\sum_{j=1}^{t}\beta_j\x_j,
 \qquad \x_j\in\F_2^k.
\end{equation*}
Because $\G$ is binary,
$\x\G=\sum_j\beta_j(\x_j\G)$.  A coordinate of $\x\G$ is zero exactly when
that coordinate is zero in every $\x_j\G$.  Hence
\begin{equation*}
 \supp(\x\G)=\bigcup_{j=1}^{t}\supp(\x_j\G).
\end{equation*}
If $\x\ne0$, some $\x_j$ is nonzero.  The union therefore has size at least
the binary minimum distance.  Conversely, binary messages form a subset of
$K^k$, which gives the reverse inequality.
\end{proof}

Indeed, the proof writes the extension-field message as $t$ binary component
messages and applies the same binary matrix to each component independently.
This decomposition is useful for binary-choice Silent OT.  It is not the
mixing required by a field-valued Silent VOLE instance.  For that application
we instead use the labeled finite-field construction of
Section~\ref{sec:prime-fields}, whose edge labels and convolution coefficients
are sampled in the ambient field.
